\documentclass[12pt, a4paper]{article}

% --- 1. Cấu hình Tiếng Việt và Font chữ ---
\usepackage[utf8]{inputenc}
\usepackage[T5]{fontenc}
\usepackage[vietnamese]{babel}
\usepackage{newtxtext,newtxmath}

\makeatletter
\renewcommand{\normalsize}{\@setfontsize\normalsize{13pt}{16pt}}
\makeatother
\normalsize

% 2. Gói Toán học & Ký hiệu bổ trợ
\usepackage{amsmath, amssymb, amsfonts, mathrsfs, amsthm}
\usepackage{graphicx}
\usepackage{fontawesome}
\usepackage{booktabs, tabularx, array, multirow}
\usepackage{enumitem}

% 3. Đồ họa TikZ, Bảng biến thiên & Hình không gian
\usepackage{tikz}
\usetikzlibrary{calc, intersections, angles, quotes, patterns, arrows.meta, 3d}
\usepackage{pgfplots}
\pgfplotsset{compat=1.18}
\usepackage{tkz-euclide}
\usepackage{tkz-tab}

\renewcommand{\vec}[1]{\overrightarrow{#1}}

% 4. Hộp sư phạm (Tcolorbox)
\usepackage[many]{tcolorbox}
\tcbuselibrary{skins, breakable}

\newtcolorbox{pedabox}[1][]{
    enhanced,
    colback=blue!5!white,
    colframe=blue!75!black,
    fonttitle=\bfseries\sffamily,
    title=#1,
    boxrule=0.8pt,
    arc=2mm,
    breakable,
    left=3mm, right=3mm, top=2mm, bottom=2mm
}

\newtcolorbox{scaffoldbox}[1][]{
    enhanced,
    colback=orange!5!white,
    colframe=orange!85!black,
    fonttitle=\bfseries\sffamily,
    title=#1,
    boxrule=0.8pt,
    arc=2mm,
    breakable,
    left=3mm, right=3mm, top=2mm, bottom=2mm
}

\newtcolorbox{aibox}[1][]{
    enhanced,
    colback=green!5!white,
    colframe=teal!80!black,
    fonttitle=\bfseries\sffamily,
    title=#1,
    boxrule=0.8pt,
    arc=2mm,
    breakable,
    left=3mm, right=3mm, top=2mm, bottom=2mm
}

% 5. Căn lề chuẩn văn bản giáo án
\usepackage{geometry}
\geometry{
    a4paper,
    left=25mm,
    right=15mm,
    top=20mm,
    bottom=20mm
}

% 6. Header / Footer chuẩn hóa (BẮT BUỘC CHÍNH XÁC NỘI DUNG NÀY)
\usepackage{fancyhdr}
\usepackage{lastpage}
\pagestyle{fancy}
\fancyhf{}

\fancyhead[L]{\small \textit{Kế hoạch bài dạy Toán 11}}
\fancyhead[R]{\textbf{Trang \thepage/\pageref{LastPage}}}
\fancyfoot[L]{\small \textit{Tổ Toán - Tin}}
\fancyfoot[R]{\small \textit{Trường THCS\&THPT Hồng Vân}}
\renewcommand{\headrulewidth}{0.4pt}
\renewcommand{\footrulewidth}{0.4pt}

% 7. Giãn dòng & Giãn đoạn theo chuẩn Nghị định 30/2020/NĐ-CP
\usepackage{setspace}
\setstretch{1.15}
\setlength{\parindent}{1.0cm}
\setlength{\parskip}{5pt}

\begin{document}

\begin{center}
    {\Large \textbf{KẾ HOẠCH BÀI DẠY MÔN TOÁN LỚP 11}}\\[0.4em]
    {\LARGE \textbf{BÀI 12. ĐƯỜNG THẲNG VÀ MẶT PHẲNG}}\\[0.3em]
    {\LARGE \textbf{SONG SONG}}\\[0.5em]
    \textbf{Môn học:} Toán 11 | \textbf{Thời lượng:} 02 tiết (Tiết 31, 32 theo PPCT)
\end{center}

\vspace{0.5em}
\hrule
\vspace{1em}

\section*{I. MỤC TIÊU DẠY HỌC}

\subsection*{1. Kiến thức, Năng lực}

\textbf{a) Năng lực Toán học đặc thù (nguyên văn chuẩn CT GDPT 2018):}
\begin{itemize}[leftmargin=1.5em]
    \item \textbf{Năng lực tư duy và lập luận toán học:}
    \begin{itemize}
        \item Nhận biết được vị trí tương đối của đường thẳng và mặt phẳng trong không gian: đường thẳng cắt mặt phẳng (có 1 điểm chung), đường thẳng nằm trong mặt phẳng (có vô số điểm chung), đường thẳng song song với mặt phẳng (không có điểm chung).
        \item Hiểu rõ và phát biểu chính xác định nghĩa đường thẳng song song với mặt phẳng ($d \parallel (\alpha) \iff d \cap (\alpha) = \varnothing$).
        \item Giải thích được điều kiện để đường thẳng song song với mặt phẳng (Định lí 1): Nếu đường thẳng $d$ không nằm trong mặt phẳng $(\alpha)$ và song song với một đường thẳng $d'$ nằm trong $(\alpha)$ thì $d \parallel (\alpha)$.
        \item Giải thích được các tính chất cơ bản của đường thẳng song song với mặt phẳng:
        \begin{itemize}
            \item Định lí 2: Cho đường thẳng $d$ song song với mặt phẳng $(\alpha)$. Nếu mặt phẳng $(\beta)$ chứa $d$ và cắt $(\alpha)$ theo giao tuyến $d'$ thì $d' \parallel d$.
            \item Hệ quả: Nếu hai mặt phẳng phân biệt cùng song song với một đường thẳng thì giao tuyến của chúng (nếu có) cũng song song với đường thẳng đó.
        \end{itemize}
    \end{itemize}
    \item \textbf{Năng lực giải quyết vấn đề toán học:}
    \begin{itemize}
        \item Nắm vững phương pháp then chốt để chứng minh đường thẳng song song với mặt phẳng: chuyển bài toán chứng minh đường thẳng song song với mặt phẳng về bài toán chứng minh hai đường thẳng song song trong không gian.
        \item Vận dụng định lí và hệ quả để tìm giao tuyến của hai mặt phẳng, xác định thiết diện của hình chóp, hình tứ diện cắt bởi mặt phẳng song song với một đường thẳng cho trước.
    \end{itemize}
    \item \textbf{Năng lực mô hình hóa toán học:}
    \begin{itemize}
        \item Vận dụng được kiến thức về đường thẳng song song với mặt phẳng để mô tả các hình ảnh thực tiễn: thanh rèm cửa song song với sàn nhà/trần nhà, xà gồ mái nhà song song với mặt sàn, dây cáp điện ngắm dọc theo tuyến đường, lan can cầu thang.
    \end{itemize}
    \item \textbf{Năng lực giao tiếp toán học:}
    \begin{itemize}
        \item Sử dụng chính xác các thuật ngữ và kí hiệu: $d \parallel (P)$, $d \cap (P) = \{A\}$, $d \subset (P)$; diễn đạt mạch lạc các bước lập luận hình học.
    \end{itemize}
    \item \textbf{Năng lực sử dụng công cụ, phương tiện học toán:}
    \begin{itemize}
        \item Sử dụng thước thẳng vẽ hình biểu diễn chính xác (giữ nguyên tính song song); sử dụng thiết bị số (smartphone) chụp và truyền chiếu hình ảnh thực tế; dùng phần mềm GeoGebra 3D để trực quan hóa quan hệ song song.
    \end{itemize}
\end{itemize}

\textbf{b) Năng lực số (theo Thông tư 02/2025/TT-BGDĐT):}
\begin{itemize}[leftmargin=1.5em]
    \item \textbf{Mã 2.2NC1a (Chia sẻ và truyền tải nội dung số qua các thiết bị kết nối không dây):} Học sinh sử dụng điện thoại thông minh (smartphone) chụp ảnh các hình ảnh thực tế về đường thẳng song song với mặt phẳng (mép bàn, xà ngang phòng học, gờ bảng) và trình chiếu không dây (qua AirPlay, Miracast, Google Cast hoặc gửi nhanh qua nhóm học tập số) lên màn hình lớn của lớp học để cả lớp cùng phân tích, kiểm chứng.
\end{itemize}

\textbf{c) Năng lực AI (theo Quyết định 2422/QĐ-BGDĐT):}
\begin{itemize}[leftmargin=1.5em]
    \item \textbf{Mã 11.C3.1 (Kĩ thuật đặt câu lệnh Prompting và khai thác trợ lí AI trong học toán):} Học sinh học cách cấu trúc câu lệnh prompt rõ ràng, chuẩn mực (gồm bối cảnh, giả thiết hình học, yêu cầu chứng minh và cấu trúc lập luận) để nhờ trợ lí AI gợi ý sơ đồ tư duy chứng minh đường thẳng song song với mặt phẳng; biết cách phân tích phản biện để phát hiện và chỉnh sửa các lỗi suy luận logic của AI.
\end{itemize}

\textbf{d) Năng lực chung \& Phẩm chất:}
\begin{itemize}[leftmargin=1.5em]
    \item \textbf{Năng lực chung:} Tự chủ và tự học (chủ động tìm kiếm đường thẳng song song trung gian); giao tiếp và hợp tác nhóm khi cùng giải bài toán thiết diện.
    \item \textbf{Phẩm chất:} Trung thực, cẩn thận trong suy luận logic hình học; trách nhiệm trong học tập.
\end{itemize}

\section*{II. THIẾT BỊ DẠY HỌC VÀ HỌC LIỆU}
\begin{itemize}[leftmargin=1.5em]
    \item \textbf{Giáo viên:} Kế hoạch bài dạy chuẩn CV 5512; bài trình chiếu PowerPoint trực quan; mô hình khung không gian (hình chóp, hình hộp); màn hình tương tác/máy chiếu kết nối không dây; Phiếu học tập số 1, 2; Rubric đánh giá năng lực.
    \item \textbf{Học sinh:} Vở ghi bài, SGK Toán 11; thước kẻ có chia độ dài; bút chì, tẩy; điện thoại thông minh (mỗi nhóm 1 thiết bị) có kết nối mạng Internet.
\end{itemize}

\section*{III. TIẾN TRÌNH DẠY HỌC (TRỌN VẸN 02 TIẾT)}

\begin{center}
    \framebox[1.0\textwidth]{\parbox{0.96\textwidth}{
        \textbf{CẤU TRÚC PHÂN PHỐI 02 TIẾT DẠY HỌC (TIẾT 31, 32 THEO PPCT):}\\[0.3em]
        $\bullet$ \textbf{Tiết 1 (Tiết 31 PPCT):} Vị trí tương đối của đường thẳng và mặt phẳng --- Điều kiện để đường thẳng song song với mặt phẳng (Định lí 1) --- Năng lực số trình chiếu ảnh thực tế (Mã 2.2NC1a).\\[0.3em]
        $\bullet$ \textbf{Tiết 2 (Tiết 32 PPCT):} Tính chất của đường thẳng song song với mặt phẳng (Định lí 2 và Hệ quả) --- Kĩ thuật Prompting AI (Mã 11.C3.1) --- Bài toán thiết diện song song --- Luyện tập tổng hợp.
    }}
\end{center}

\vspace{0.8em}
\hrule
\vspace{0.8em}

{\large \textbf{TIẾT 1 (TIẾT 31 THEO PHÂN PHỐI CHƯƠNG TRÌNH)}}\\[0.3em]
\textbf{Nội dung:} Vị trí tương đối của đường thẳng và mặt phẳng --- Điều kiện để đường thẳng song song với mặt phẳng --- Năng lực số (Mã 2.2NC1a).

\subsubsection*{1. Hoạt động 1: Mở đầu / Khởi động (05 phút)}
\textbf{Chủ đề:} \textit{Cây gác cờ và thanh rèm cửa sổ trong phòng học}
\begin{itemize}
    \item \textbf{a) Mục tiêu:} Tạo tình huống trực quan trong đời sống lớp học để học sinh nhận thức được khái niệm đường thẳng không có điểm chung với một mặt phẳng (đường thẳng song song với mặt phẳng).
    \item \textbf{b) Nội dung:} Giáo viên yêu cầu học sinh quan sát thanh treo rèm cửa sổ phía trên tường:
    \begin{itemize}
        \item Thanh treo rèm có cắt mặt phẳng sàn nhà không?
        \item Thanh treo rèm có nằm trên mặt sàn nhà không?
        \item Quan hệ vị trí giữa thanh treo rèm và mặt sàn nhà là gì?
    \end{itemize}
    \item \textbf{c) Sản phẩm:} Học sinh trả lời: Thanh treo rèm không bao giờ chạm tới sàn nhà, tức là không có điểm chung với sàn nhà. Ta nói thanh treo rèm song song với mặt sàn.
    \item \textbf{d) Tổ chức thực hiện:} GV đặt câu hỏi $\to$ HS quan sát, trao đổi nhanh $\to$ GV dẫn dắt vào bài học mới.
\end{itemize}

\subsubsection*{2. Hoạt động 2: Hình thành kiến thức mới (25 phút)}

\paragraph{Mục 1: Vị trí tương đối của đường thẳng và mặt phẳng (10 phút)}
Cho đường thẳng $d$ và mặt phẳng $(\alpha)$. Dựa vào số điểm chung, có ba vị trí tương đối:
\begin{enumerate}[label=\textbf{\arabic*.}]
    \item \textbf{Đường thẳng cắt mặt phẳng:} $d$ và $(\alpha)$ có \textbf{duy nhất một điểm chung} $A$.
    Kí hiệu: $d \cap (\alpha) = \{A\}$ hoặc $d \cap (\alpha) = A$.
    \item \textbf{Đường thẳng nằm trong mặt phẳng:} $d$ và $(\alpha)$ có \textbf{vô số điểm chung} (mọi điểm của $d$ đều thuộc $(\alpha)$).
    Kí hiệu: $d \subset (\alpha)$ hoặc $(\alpha) \supset d$.
    \item \textbf{Đường thẳng song song với mặt phẳng:} $d$ và $(\alpha)$ \textbf{không có điểm chung nào}.
    Kí hiệu: $d \parallel (\alpha)$ hoặc $(\alpha) \parallel d$.
    $$\text{Định nghĩa: } d \parallel (\alpha) \iff d \cap (\alpha) = \varnothing.$$
\end{enumerate}

\begin{center}
\begin{tikzpicture}[scale=0.85]
    % Cắt nhau
    \draw[thick] (-0.2,-0.2) -- (2.8,-0.2) -- (3.8,1.8) -- (0.8,1.8) -- cycle;
    \node at (0.8,1.4) {$(\alpha)$};
    \draw[thick, blue] (1.8, 2.5) -- (1.8, 0.8);
    \draw[dashed, thick, blue] (1.8, 0.8) -- (1.8, 0.2);
    \draw[thick, blue] (1.8, 0.2) -- (1.8, -0.6) node[below] {$d$};
    \fill[red] (1.8, 0.8) circle (2pt) node[right=2pt] {$A$};
    \node at (1.8,-1.2) {\textbf{a) $d$ cắt $(\alpha)$ tại $A$}};

    % Nằm trong
    \begin{scope}[xshift=5.2cm]
    \draw[thick] (-0.2,-0.2) -- (2.8,-0.2) -- (3.8,1.8) -- (0.8,1.8) -- cycle;
    \node at (0.8,1.4) {$(\alpha)$};
    \draw[thick, blue] (0.3, 0.6) -- (3.3, 1.0) node[right] {$d$};
    \node at (1.8,-1.2) {\textbf{b) $d$ nằm trong $(\alpha)$ ($d \subset (\alpha)$)}};
    \end{scope}

    % Song song
    \begin{scope}[xshift=10.4cm]
    \draw[thick] (-0.2,-0.2) -- (2.8,-0.2) -- (3.8,1.8) -- (0.8,1.8) -- cycle;
    \node at (0.8,1.4) {$(\alpha)$};
    \draw[thick, blue] (0.3, 2.3) -- (3.3, 2.3) node[right] {$d$};
    \node at (1.8,-1.2) {\textbf{c) $d$ song song $(\alpha)$ ($d \parallel (\alpha)$)}};
    \end{scope}
\end{tikzpicture}
\end{center}

\paragraph{Mục 2: Điều kiện để đường thẳng song song với mặt phẳng (Định lí 1) (10 phút)}
\begin{itemize}
    \item \textbf{Định lí 1 (Dấu hiệu nhận biết):} Nếu đường thẳng $d$ không nằm trong mặt phẳng $(\alpha)$ và song song với một đường thẳng $d'$ nằm trong $(\alpha)$ thì $d$ song song với $(\alpha)$.
    $$\begin{cases} d \not\subset (\alpha) \\ d' \subset (\alpha) \\ d \parallel d' \end{cases} \implies d \parallel (\alpha).$$
\end{itemize}

\begin{scaffoldbox}[Giàn giáo sư phạm cho HS TB - Yếu: Phương pháp 3 bước chứng minh $d \parallel (\alpha)$]
    Để chứng minh đường thẳng $d$ song song với mặt phẳng $(\alpha)$, học sinh thực hiện theo 3 bước sau:\\
    $\bullet$ \textbf{Bước 1 (Kiểm tra điều kiện ngoài):} Khẳng định $d$ không nằm trong mặt phẳng $(\alpha)$ ($d \not\subset (\alpha)$).\\
    $\bullet$ \textbf{Bước 2 (Tìm đường thẳng song song trong mặt phẳng):} Tìm trong $(\alpha)$ một đường thẳng $d'$ sao cho $d \parallel d'$ (dùng đường trung bình của tam giác, hình bình hành, hoặc tính chất bắc cầu).\\
    $\bullet$ \textbf{Bước 3 (Kết luận):} Áp dụng Định lí 1 suy ra $d \parallel (\alpha)$.\\
    \textit{Khắc sâu:} Muốn chứng minh ĐƯỜNG song song với MẶT, ta chỉ cần chứng minh ĐƯỜNG song song với một ĐƯỜNG nằm trong MẶT!
\end{scaffoldbox}

\paragraph{Mục 3: Tích hợp Năng lực số (Mã 2.2NC1a) (05 phút)}

\begin{pedabox}[Năng lực số (Mã 2.2NC1a): Chụp và trình chiếu hình ảnh đường song song mặt phẳng]
    $\bullet$ \textbf{Nhiệm vụ theo nhóm (03 phút):}\\
    1. Đại diện mỗi nhóm dùng smartphone chụp một hình ảnh thực tế trong lớp học thể hiện đường thẳng song song với một mặt phẳng (ví dụ: gờ đèn tuýp song song mặt sàn, cạnh bàn học song song với bảng đen, mép cửa sổ song song với tường đối diện).\\
    2. Sử dụng tính năng phản chiếu màn hình (Screen Mirroring / AirPlay / Google Cast) hoặc gửi ảnh vào nhóm Zalo/Google Classroom của lớp để trình chiếu lên màn hình tivi/máy chiếu.\\
    $\bullet$ \textbf{Thuyết trình nhanh (02 phút):}\\
    Nhóm chỉ rõ: Đâu là đường thẳng $d$? Đâu là mặt phẳng $(\alpha)$? Đâu là đường thẳng $d' \subset (\alpha)$ đóng vai trò song song để áp dụng Định lí 1?
\end{pedabox}

\subsubsection*{3. Hoạt động 3: Luyện tập (10 phút)}
\begin{itemize}
    \item \textbf{Bài toán 1:} Cho hình chóp $S.ABCD$ có đáy $ABCD$ là hình bình hành. Gọi $M, N$ lần lượt là trung điểm của các cạnh $SA$ và $SB$.
    \begin{itemize}
        \item a) Chứng minh rằng $MN \parallel (ABCD)$.
        \item b) Chứng minh rằng $AB \parallel (SCD)$ và $CD \parallel (SAB)$.
    \end{itemize}
    \item \textit{Lời giải chi tiết:}\\
    a) Chứng minh $MN \parallel (ABCD)$:\\
    - Trong tam giác $SAB$, $M$ là trung điểm $SA$, $N$ là trung điểm $SB \implies MN$ là đường trung bình của tam giác $SAB$.\\
    Suy ra: $MN \parallel AB$.\\
    - Ta có: $\begin{cases} MN \not\subset (ABCD) \\ AB \subset (ABCD) \\ MN \parallel AB \end{cases} \implies MN \parallel (ABCD)$ (theo Định lí 1).\\
    b) Chứng minh $AB \parallel (SCD)$:\\
    - Vì $ABCD$ là hình bình hành nên $AB \parallel CD$.\\
    - Ta có: $\begin{cases} AB \not\subset (SCD) \\ CD \subset (SCD) \\ AB \parallel CD \end{cases} \implies AB \parallel (SCD)$.\\
    - Tương tự, vì $CD \parallel AB$ mà $AB \subset (SAB)$ và $CD \not\subset (SAB) \implies CD \parallel (SAB)$.
\end{itemize}

\subsubsection*{4. Hoạt động 4: Vận dụng (05 phút)}
\textbf{Ứng dụng trong xây dựng và thi công trần nhà:}\\
Khi đóng trần thạch cao, những người thợ bắn các thanh xương cá kim loại treo trên cao. Để bảo đảm toàn bộ hệ trần phẳng song song với sàn nhà, họ căng các sợi dây thừng cân mực nước song song với các mép tường sàn. Học sinh giải thích cơ sở toán học của việc cân mực này theo Định lí 1.

\newpage

{\large \textbf{TIẾT 2 (TIẾT 32 THEO PHÂN PHỐI CHƯƠNG TRÌNH)}}\\[0.3em]
\textbf{Nội dung:} Tính chất của đường thẳng song song với mặt phẳng (Định lí 2 và Hệ quả) --- Năng lực AI (Mã 11.C3.1) --- Thiết diện song song --- Luyện tập tổng hợp.

\subsubsection*{1. Hoạt động 1: Mở đầu / Khởi động (05 phút)}
\textbf{Chủ đề:} \textit{Vết cắt của lưỡi cưa máy trên thân gỗ tròn phẳng}
\begin{itemize}
    \item \textbf{a) Mục tiêu:} Giúp học sinh dự đoán được mối quan hệ giữa đường thẳng song song có sẵn với giao tuyến của mặt phẳng cắt.
    \item \textbf{b) Nội dung:} Cho đường thẳng $d$ song song với mặt phẳng bàn $(\alpha)$. Một tấm bìa cứng $(\beta)$ xoay quanh đường thẳng $d$ và cắt mặt bàn $(\alpha)$ theo giao tuyến $d'$. Giao tuyến $d'$ này có vị trí tương đối như thế nào với đường thẳng $d$?
    \item \textbf{c) Sản phẩm:} Học sinh dự đoán: Giao tuyến $d'$ phải song song với $d$.
    \item \textbf{d) Tổ chức thực hiện:} GV thao tác mô phỏng $\to$ HS phát biểu $\to$ GV dẫn dắt vào Định lí 2 (tính chất cốt lõi để tìm giao tuyến và thiết diện).
\end{itemize}

\subsubsection*{2. Hoạt động 2: Hình thành kiến thức mới (18 phút)}

\paragraph{Mục 1: Tính chất của đường thẳng song song với mặt phẳng (Định lí 2) (06 phút)}
\begin{itemize}
    \item \textbf{Định lí 2:} Cho đường thẳng $d$ song song với mặt phẳng $(\alpha)$. Nếu mặt phẳng $(\beta)$ chứa $d$ và cắt $(\alpha)$ theo giao tuyến $d'$ thì $d'$ song song với $d$.
    $$\begin{cases} d \parallel (\alpha) \\ d \subset (\beta) \\ (\alpha) \cap (\beta) = d' \end{cases} \implies d' \parallel d.$$
    \item \textbf{Hệ quả (Giao tuyến của hai mặt phẳng cùng song song với một đường thẳng):} Nếu hai mặt phẳng phân biệt cùng song song với một đường thẳng thì giao tuyến của chúng (nếu có) cũng song song với đường thẳng đó.
    $$\begin{cases} (\alpha) \parallel d \\ (\beta) \parallel d \\ (\alpha) \cap (\beta) = d' \end{cases} \implies d' \parallel d.$$
\end{itemize}

\begin{center}
\begin{tikzpicture}[scale=0.9]
    % Mặt phẳng (alpha)
    \draw[thick] (-0.2,-0.2) -- (4,-0.2) -- (5.2,1.8) -- (1,1.8) -- cycle;
    \node at (1.2,1.4) {$(\alpha)$};

    % Đường d song song
    \draw[thick, blue] (0.5, 3.2) -- (4.5, 3.2) node[right] {$d$};

    % Mặt phẳng (beta) chứa d cắt (alpha)
    \draw[thick, fill=blue!10, opacity=0.4] (0.5, 3.2) -- (4.5, 3.2) -- (4.0, 0.8) -- (0.0, 0.8) -- cycle;
    \draw[thick] (0.5, 3.2) -- (0.0, 0.8);
    \draw[thick] (4.5, 3.2) -- (4.0, 0.8);
    \node at (0.3, 2.5) {$(\beta)$};

    % Giao tuyến d'
    \draw[thick, red] (0.0, 0.8) -- (4.0, 0.8) node[right] {$d'$};

    \node at (2.5, -0.8) {\textbf{Định lí 2: $d \parallel (\alpha), d \subset (\beta) \implies d' \parallel d$}};
\end{tikzpicture}
\end{center}

\paragraph{Mục 2: Tích hợp Năng lực AI (Mã 11.C3.1) (06 phút)}

\begin{aibox}[Tích hợp Năng lực AI (QĐ 2422/QĐ-BGDĐT --- Mã 11.C3.1): Kĩ thuật Prompting AI hỗ trợ giải toán]
    $\bullet$ \textbf{Tình huống học tập:}\\
    Khi gặp một bài toán hình học không gian phức tạp về chứng minh đường thẳng song song với mặt phẳng, làm thế nào để khai thác trợ lí AI hiệu quả mà không bị phụ thuộc?\\
    $\bullet$ \textbf{Cấu trúc Prompt chuẩn mực (Học sinh thực hành trên máy tính/điện thoại):}\\
    \textit{``Tôi đang học Toán 11 bài Đường thẳng và mặt phẳng song song. Cho hình chóp $S.ABCD$ đáy hình bình hành, $M$ là trung điểm $SC$. Hãy hướng dẫn tôi sơ đồ tư duy (bước 1, bước 2) để chứng minh $SA \parallel (BDM)$ mà không giải tắt. Giải thích rõ vì sao phải gọi tâm đáy $O$.''}\\
    $\bullet$ \textbf{Đánh giá kết quả của AI:}\\
    Học sinh kiểm tra: AI có nhận ra $OM$ là đường trung bình của tam giác $SAC$ hay không? Nhận biết giới hạn của AI (đôi khi AI có thể nhầm lẫn vị trí điểm nếu không mô tả chính xác giả thiết hình học).
\end{aibox}

\paragraph{Mục 3: Phương pháp xác định thiết diện song song với một đường thẳng (06 phút)}
\begin{itemize}
    \item \textbf{Nhiệm vụ:} Xác định thiết diện của hình chóp cắt bởi mặt phẳng $(\alpha)$ đi qua điểm $M$ và song song với đường thẳng $d$.
    \item \textbf{Phương pháp giải:}
    \begin{itemize}
        \item Tìm mặt phẳng của hình chóp chứa điểm $M$ và cắt đường thẳng $d$ (hoặc chứa $d$).
        \item Áp dụng Định lí 2: Qua $M$, kẻ đường thẳng song song với $d$, đường này cắt các cạnh của mặt đó tại các đỉnh mới của thiết diện.
        \item Tiếp tục mở rộng sang các mặt kề bên cho đến khi tạo thành đa giác giao tuyến khép kín.
    \end{itemize}
\end{itemize}

\subsubsection*{3. Hoạt động 3: Luyện tập (16 phút)}
\begin{itemize}
    \item \textbf{Bài toán trọng tâm:} Cho hình chóp $S.ABCD$ có đáy $ABCD$ là hình bình hành tâm $O$. Gọi $M$ là trung điểm của cạnh $SC$.
    \begin{itemize}
        \item a) Chứng minh rằng $OM \parallel (SAB)$ và $OM \parallel (SAD)$.
        \item b) Gọi $(\alpha)$ là mặt phẳng đi qua $M$ và song song với hai đường thẳng $BD$ và $SA$. Xác định thiết diện của hình chóp cắt bởi mặt phẳng $(\alpha)$. Thiết diện là hình gì?
    \end{itemize}
    \item \textit{Lời giải chi tiết:}\\
    a) Chứng minh $OM \parallel (SAB)$ và $OM \parallel (SAD)$:\\
    - Trong tam giác $SAC$, vì $ABCD$ là hình bình hành nên $O$ là trung điểm của $AC$. Mặt khác $M$ là trung điểm của $SC$.\\
    Suy ra: $OM$ là đường trung bình của tam giác $SAC \implies OM \parallel SA$.\\
    - Ta có: $\begin{cases} OM \not\subset (SAB) \\ SA \subset (SAB) \\ OM \parallel SA \end{cases} \implies OM \parallel (SAB)$.\\
    - Tương tự, vì $SA \subset (SAD)$ nên $OM \parallel (SAD)$.\\
    b) Xác định thiết diện cắt bởi $(\alpha)$:\\
    - Điểm $M \in SC$ thuộc mặt phẳng $(\alpha)$.\\
    - Vì $(\alpha) \parallel BD$ và $(\alpha) \parallel SA$:\\
      + Trong mặt phẳng $(SBD)$, mặt phẳng $(\alpha)$ chứa đường thẳng đi qua $M$ là không tiện, ta xét các mặt phẳng chứa điểm $M$:\\
      + Mặt phẳng $(SCD)$ chứa $M$: Trong $(ABCD)$ có $BD \parallel (\alpha)$, trong $(SCD)$ chưa có đường song song, ta xét mặt phẳng $(SBC)$:\\
      Trong $(SBC)$: Qua $M$ kẻ đường thẳng song song với $SA$ không được (vì $SA \not\subset (SBC)$).\\
      + Ta xét mặt phẳng phụ $(SAC)$ chứa $M$: Qua $M$ kẻ đường thẳng song song với $SA$ cắt $AC$ tại $H$. Vì $M$ là trung điểm $SC$ nên $H$ là trung điểm $OC$.\\
      + Qua $H$, kẻ đường thẳng song song với $BD$ cắt $BC$ tại $P$ và cắt $CD$ tại $Q$.\\
      + Trong $(SBC)$: Qua $P$ kẻ đường thẳng song song với $SA$ cắt $SB$ tại $E$.\\
      + Trong $(SCD)$: Qua $Q$ kẻ đường thẳng song song với $SA$ cắt $SD$ tại $F$.\\
    - Nối các đoạn giao tuyến: $PQ \subset (ABCD)$, $QE \dots$ Cụ thể, các đoạn giao tuyến cắt các mặt là: $P, Q, F, M, E$.\\
    Đoạn giao tuyến là ngũ giác $EPQFM$ (hoặc hình thang/hình bình hành tùy theo vị trí cắt).\\
    Khi xét cắt bởi mặt phẳng đi qua $M$ song song với $BD$ và $SA$, thiết diện thu được là một ngũ giác có các cạnh đối diện song song từng đôi một với $BD$ và $SA$.
\end{itemize}

\subsubsection*{4. Hoạt động 4: Vận dụng (05 phút)}
\textbf{Thiết kế máng xối thoát nước mái tôn:}\\
Máng xối thoát nước tôn được đặt dọc theo mép mái nhà và song song với mặt sàn nhà. Dựa vào tính chất giao tuyến song song (Định lí 2), học sinh giải thích cách người thợ lắp đặt máng xối sao cho dòng chảy luôn thẳng hướng, không bị ứ đọng nước hoặc lệch độ dốc so với mép tường bao.

\newpage

\section*{IV. PHỤ LỤC HỌC LIỆU BỔ TRỢ}

\subsection*{Phụ lục 1: Phiếu học tập số 1 (Tiết 31) --- Vị trí tương đối \& Dấu hiệu song song}
\noindent\textbf{Họ và tên:} .................................................................... \textbf{Lớp:} 11A....... \textbf{Nhóm:} ..........

\begin{pedabox}[Nhiệm vụ 1: Hoàn thành bảng vị trí tương đối của đường thẳng và mặt phẳng]
\begin{center}
\begin{tabularx}{0.95\linewidth}{|p{4.2cm}|X|X|}
\hline
\textbf{Vị trí tương đối} & \textbf{Số điểm chung} & \textbf{Kí hiệu toán học} \\
\hline
Đường thẳng cắt mặt phẳng & Có \dots\dots\dots\dots\dots\dots điểm chung & $d \cap (\alpha) = \{A\}$ \\
\hline
Đường thẳng nằm trong mặt phẳng & Có \dots\dots\dots\dots\dots\dots điểm chung & $d \subset (\alpha)$ \\
\hline
Đường thẳng song song mặt phẳng & Có \dots\dots\dots\dots\dots\dots điểm chung & $d \parallel (\alpha)$ \\
\hline
\end{tabularx}
\end{center}
\end{pedabox}

\begin{pedabox}[Nhiệm vụ 2: Điền sơ đồ 3 bước chứng minh $d \parallel (\alpha)$]
$\bullet$ Bước 1: Chứng minh $d$ không nằm trong $(\alpha)$ ($d \not\subset \dots\dots\dots$).\\
$\bullet$ Bước 2: Tìm trong $(\alpha)$ một đường thẳng $d'$ sao cho $d \parallel \dots\dots\dots$\\
$\bullet$ Bước 3: Kết luận $d \parallel (\alpha)$ theo Định lí \dots\dots\dots
\end{pedabox}

\subsection*{Phụ lục 2: Phiếu học tập số 2 (Tiết 32) --- Định lí 2 và Tìm thiết diện}
\noindent\textbf{Họ và tên:} .................................................................... \textbf{Lớp:} 11A....... \textbf{Nhóm:} ..........

\begin{pedabox}[Nhiệm vụ: Áp dụng Định lí 2 hoàn thành phát biểu sau]
$\bullet$ Cho đường thẳng $d \parallel (\alpha)$. Mặt phẳng $(\beta)$ chứa $d$ và cắt $(\alpha)$ theo giao tuyến $d'$.\\
Khi đó: $d'$ \dots\dots\dots\dots\dots\dots\dots\dots với đường thẳng $d$.\\[0.3em]
$\bullet$ \textbf{Bài tập áp dụng:} Cho hình chóp $S.ABCD$ có đáy là hình bình hành. Gọi $(\alpha)$ là mặt phẳng qua trung điểm $M$ của $AB$ và song song với $BC$.\\
Mặt phẳng $(\alpha)$ cắt $CD$ tại $N$. Đoạn thẳng $MN$ có vị trí tương đối như thế nào với $BC$? Vì sao?
\end{pedabox}

\subsection*{Phụ lục 3: Bảng Rubric đánh giá năng lực tích hợp trong Bài 12}

\begin{center}
\begin{tabularx}{\linewidth}{|p{3.2cm}|X|X|X|X|}
\hline
\textbf{Tiêu chí} & \textbf{Mức 1 (Chưa đạt)} & \textbf{Mức 2 (Đạt)} & \textbf{Mức 3 (Khá)} & \textbf{Mức 4 (Tốt)} \\
\hline
\textbf{1. Nhận biết và phát biểu định lí} & Nhầm lẫn giữa đường thẳng song song mặt phẳng với hai đường thẳng song song. & Phát biểu đúng định nghĩa và Định lí 1, Định lí 2. & Giải thích rõ ràng điều kiện không chứa $d \not\subset (\alpha)$ trong lập luận. & Vận dụng linh hoạt hệ quả và phân tích sâu sắc các trường hợp suy biến. \\
\hline
\textbf{2. Kĩ năng chứng minh \& Dựng thiết diện} & Chưa biết tìm đường thẳng song song trung gian trong mặt phẳng. & Chứng minh được các bài toán cơ bản dùng đường trung bình tam giác. & Thành thạo phương pháp chứng minh $d \parallel (P)$ và dựng thiết diện song song cơ bản. & Dựng thiết diện phức tạp (ngũ giác, hình thang), biện luận hình dạng hình học chính xác. \\
\hline
\textbf{3. Năng lực số \& AI (Mã 2.2NC1a \& 11.C3.1)} & Chưa biết chia sẻ hình ảnh qua thiết bị số; đặt prompt AI chung chung. & Chụp và truyền chiếu được ảnh thực tế lên màn hình lớp học. & Đặt prompt AI có cấu trúc rõ ràng và kiểm chứng được các bước gợi ý của AI. & Đánh giá phản biện lời giải của AI, phát hiện lỗ hổng logic và tối ưu hóa câu lệnh prompt. \\
\hline
\end{tabularx}
\end{center}

\end{document}
